%Chargement des paquets
\usepackage{amsmath}
\usepackage{amsfonts}
\usepackage{amssymb}
\usepackage{enumerate}
\usepackage{mathtools}
\usepackage{bbm}
\usepackage{xparse, etoolbox}
\usepackage{enumerate}
\usepackage{enumitem}
\usepackage{mathabx}
\usepackage{babel}[french]

%environment
\usepackage{amsthm}
\usepackage{keytheorems}
\usepackage{hyperref} 

%Interface théorème
\newkeytheorem{lem}[
	name = Lemme
]

\newkeytheorem{prop}[
	name = Proposition
]

\newkeytheorem{thm}[
	name = Théorème
]

\newkeytheorem{coro}[
	name = Corollaire
]

\renewcommand*{\proofname}{Démonstration}

\theoremstyle{definition}
\newtheorem{defn}{Définition}

\theoremstyle{definition}
\newtheorem*{nota}{Notation}

\theoremstyle{definition}
\newtheorem*{limi}{Liminaire}

\theoremstyle{remark}
\newtheorem{rmq}{Remarque}

\theoremstyle{plain}
\newtheorem*{fait}{Fait}


%Ensembles de nombres
\newcommand{\N} {\mathbb{N}}
\newcommand{\Ne}{\N^\ast}
\newcommand{\Z} {\mathbb{Z}}
\newcommand{\Q} {\mathbb{Q}}
\newcommand{\R} {\mathbb{R}}
\newcommand{\Rb}{\overline{\mathbb{R}}}
\newcommand{\Rp}{\R_+}
\newcommand{\Rm}{\R_-}
\newcommand{\K} {\mathbb{K}}
\newcommand{\Ket} {\mathbb{K^{\ast}}}
\newcommand{\Cx}{\mathbb{C}}

%Ensembles, tribus
\newcommand{\A}{\mathbb{A}}
\renewcommand{\P}{\mathcal{P}}
\newcommand{\T}{\mathcal{T}}
\newcommand{\F}{\mathcal{F}}
\newcommand{\C} {\mathcal{C}}
\newcommand{\ouv}{\mathcal{O}}
\newcommand{\B}{\mathcal{B}}
\newcommand{\M}{\mathcal{M}}
\newcommand{\etag}{\mathcal{E}}
\newcommand{\etagOp}{\etag^{+}(\Omega)}
\newcommand{\etagO}{\etag(\Omega)}
%%Lp avant quotientage
\newcommand{\Lic}{\mathcal{L}^1}
\newcommand{\Lpc}{\mathcal{L}^p}
\newcommand{\Linfc}{\mathcal{L}^{\infty}}
\newcommand{\Lifull}{\Lic(\Omega, \B, m)}
\newcommand{\Lpfull}{\Lpc(\Omega, \B, m)}
\newcommand{\Linffull}{\Linfc(\Omega, \B, m)}
%%Lp après quotientage
\newcommand{\Li}{L^1}
\newcommand{\Ld}{L^2}
\newcommand{\Lp}{L^p}
\newcommand{\Linf}{L^{\infty}}
\newcommand{\Listd}{\Li(\Omega, m)} 
\newcommand{\Ldstd}{\Ld(\Omega, m)} 
\newcommand{\Lpstd}{\Lp(\Omega, m)} 

%Quelques opérateurs
\DeclareMathOperator{\codim}{codim}
\DeclareMathOperator{\rg}{rg}
\DeclareMathOperator{\tr}{tr}
\DeclareMathOperator{\vect}{Vect}
\DeclareMathOperator{\ima}{Im}
\DeclareMathOperator{\id}{Id}
\DeclareMathOperator{\sign}{sign}
\DeclareMathOperator{\ext}{ext}
\newcommand{\ind}{\mathbbm{1}}
\newcommand*{\spr}[2]{\left(\,#1\,\middle|\,#2\,\right)}
\newcommand{\interior}[1]{%
  {\kern0pt#1}^{\mathrm{o}}%
}
\DeclareMathOperator{\card}{card}
\DeclareMathOperator{\ppcm}{ppcm}

%Commandes de pure flemme
\newcommand{\veps}{\varepsilon}
\newcommand{\intab}{\int_{a}^{b}}
\newcommand{\intR}{\int_{-\infty}^{\infty}}
\newcommand{\intinf}{\int_{- \infty}^{+ \infty}}
\newcommand{\equi}{\Leftrightarrow}
\newcommand{\Kev}{$\K$-espace vectoriel }
\newcommand{\Kevs}{$\K$-espaces vectoriels }
\newcommand{\Rev}{$\R$-espace vectoriel }
\newcommand{\Revs}{$\R$-espaces vectoriels }
\newcommand{\Cev}{$\Cx$-espace vectoriel }
\newcommand{\Cevs}{$\Cx$-espaces vectoriels }
\newcommand{\evn}{(E, \norm{.})}
\newcommand{\interf}[2]{[#1,#2]}
\newcommand{\limb}{\lim\limits_}
\newcommand{\lebn}{\lambda_n}
\newcommand{\pp}{\text{~presque partout}}
\newcommand{\derivp}[2]{\frac{\partial #1}{\partial #2}}
\newcommand{\famiu}{(u_i)_{i \in I}}
\newcommand{\sev}{\text{~s.e.v.}}
\newcommand{\Mnp}{M_{n,p}(\K)}
\newcommand{\Mn}{M_{n}(\K)}
\newcommand{\tq}{\text{~t.q.~}}
\newcommand{\mq}{~montrer~que~}
\newcommand{\et}{\quad \text{et} \quad}
\newcommand{\ou}{\text{~ou~}}
\newcommand{\Kx}{\K[X]}


%Commandes de gars chelou
\renewcommand{\subset}{\subseteq}

%Commande restriction, ty duckduckgo
\def\restr#1#2{\mathchoice
              {\setbox1\hbox{${\displaystyle #1}_{\scriptstyle #2}$}
              \restraux{#1}{#2}}
              {\setbox1\hbox{${\textstyle #1}_{\scriptstyle #2}$}
              \restraux{#1}{#2}}
              {\setbox1\hbox{${\scriptstyle #1}_{\scriptscriptstyle #2}$}
              \restraux{#1}{#2}}
              {\setbox1\hbox{${\scriptscriptstyle #1}_{\scriptscriptstyle #2}$}
              \restraux{#1}{#2}}}
\def\restraux#1#2{{#1\,\smash{\vrule height .8\ht1 depth .85\dp1}}_{\,#2}} 

%Quelques normes
\DeclarePairedDelimiter\abs{\lvert}{\rvert}
\DeclarePairedDelimiter\labs{\bigl\lvert}{\bigr\rvert}
\DeclarePairedDelimiter\norm{\lVert}{\rVert}
\DeclarePairedDelimiterXPP\normi[1]{}\lVert\rVert{_1}{\ifblank{#1}{\:\cdot\:}{#1}}
\DeclarePairedDelimiterXPP\normd[1]{}\lVert\rVert{_2}{\ifblank{#1}{\:\cdot\:}{#1}}
\DeclarePairedDelimiterXPP\normp[1]{}\lVert\rVert{_p}{\ifblank{#1}{\:\cdot\:}{#1}}
\DeclarePairedDelimiterXPP\normq[1]{}\lVert\rVert{_q}{\ifblank{#1}{\:\cdot\:}{#1}}
\DeclarePairedDelimiterXPP\norminf[1]{}\lVert\rVert{_\infty}{\ifblank{#1}{\:\cdot\:}{#1}}

%Bracket en angle
\DeclarePairedDelimiter\angl{\langle}{\rangle}

%Set, parenthèses
\newcommand{\set}[1]{\{ #1 \}}
\newcommand{\bigset}[1]{\bigl\{ #1 \bigr\}}
\newcommand{\Bigset}[1]{\Bigl\{ #1 \Bigr\}}
\newcommand{\bigpar}[1]{\bigl( #1 \bigr)}
\newcommand{\Bigpar}[1]{\Bigl( #1 \Bigr)}

%Opitmisation des opérateurs de comparaison
\DeclarePairedDelimiter\tio{o (}{)}
\DeclarePairedDelimiter\gro{O (}{)}


\pagestyle{fancy}

\fancyhead[C]{}
\fancyhead[R]{}

\begin{document}

\underline{Source :} 
\begin{itemize}
\item Gourdon - Analyse - pages 50, 57 (pour les démonstrations) ;
\item Hassan - Topologie - pages 121, 126, 128 (pour les prérequis).
\end{itemize}

\underline{Leçons :}
\begin{itemize}
\item  203 : Utilisation de la notion de compacité.
\item  206 : Exemples d’utilisation de la notion de dimension finie en analyse.
\item  208 : Espaces vectoriels normés, applications linéaires continues. Exemples.
\end{itemize}

\begin{nota}
$E$ est un \Kev et $\K$ est $\R$ ou $\Cx$.
\end{nota}

\section{Quelques prérequis de topologie}

\begin{defn}[Précompacité]
Un espace métrique $E$ est dit précompact si, pour tout $\veps>0$, il existe un recouvrement fini de $E$ par des boules (ouvertes)
de rayon $\veps$.
\end{defn}

\begin{rmq}
En particulier, tout espace compact est précompact.
\end{rmq}

\begin{lem}
Les compacts de $\R$ sont les fermés bornés de $\R$.
\end{lem}

\begin{lem}
Soient $X, Y$ deux espaces topologiques avec $Y$ séparé et $f : X \to Y$ une application continue. 
Alors, l'image par $f$ de toute partie compacte de $X$ est une partie compacte de $Y$.
\end{lem}

\begin{lem}
Soient $X, Y$ deux espaces topologiques non vides. Alors l'espace topologique produit $X \times Y$ est compact si, 
et seulement si, $X$ et $Y$ le sont.
\end{lem}

\begin{lem}
Les compacts de $\R^n$ sont les fermés bornés de $\R^n$.
\end{lem}

\begin{thm}[Equivalence des normes, page 50]
Si $E$ est de dimension finie, alors toutes les normes sur $E$ sont équivalentes.
\end{thm}

\begin{proof}
Notons $n < \infty$ la dimension de $E$ et munissons ce dernier d'une base $B=(e_1, \dots, e_n)$.
On peut donc écrire tout $x \in E$ comme $x = \sum_{i=1}^n x_i e_i$. 
L'application $N$ définie sur $E$ par $N(x) = \sup_{i \in \set{1, \dots, n}} \abs{x_i}$ est une norme sur $E$.
\begin{enumerate}[label=(\roman*)]
\item $N(x) = 0 \Rightarrow \forall i \in \set{1, \dots, n}, x_i=0 \Rightarrow x = 0$ ;
\item $\forall \lambda \in \K, \quad N(\lambda x) = \abs{\lambda} \sup_{i \in \set{1, \dots, n}} \abs{x_i} = \abs{\lambda} N(x)$ ;
\item $\forall x,y \in E, \forall i \in \set{1, \dots, n}, \quad \abs{x_i+y_i} \leq \abs{x_i} + \abs{y_i} \Rightarrow N(x+y) \leq N(x) + N(y)$.
\end{enumerate}
Soit $N'$ une seconde norme sur $E$. On veut montrer que $N'$ est équivalente à $N$.
On pose $a = \sum_i N'(e_i)$, on a alors, pour $x \in E$ :
\begin{equation} N'(x) \leq \sum_{i=1}^n N'(x_i e_i) = \sum_{i=1}^n \abs{x_i} N'(e_i) \leq a N(x) . \label{eq:1} \end{equation}
Munissons $K^n$ de la norme produit $\norminf{\cdot}$ définie par $\forall x=(x_1, \dots, x_n) \in K^n, \norminf{x} = \sup_i(x_i)$. 
C'est bien une norme sur $K^n$ (pour les mêmes raisons que $N$ est une norme sur $E$). 
L'application :
\[ \begin{array}{c c c c} 
\varphi : & (\K^n, \norminf{\cdot}) & \to & (E, N) \\
& (x_1, \dots, x_n) & \mapsto & x= \sum_{i=1}^n x_i e_i 
\end{array} , \]
est clairement une isométrie. Par ailleurs, l'ensemble $S=\set{x \in E; N(x) =1}$ est l'image de la sphère unité de $K^n$ par $\varphi$.
L'application $\varphi$ est continue (car isométrique) donc $S$ est un compact de $E$. \\
En remarquant que :
\[ \abs{N'(x) - N'(y)} \leq N'(x-y) \leq a N(x-y) \qquad \text{ d'après \ref{eq:1}} , \]
on déduit que l'application $N' : (E, N) \to \R$ est continue. Ainsi, $N'$ atteint sa borne inférieure sur le compact $S$ et comme
$x \in S \Rightarrow x \neq 0$ on en déduit que : 
\[ b = \inf_{x \in S} N'(x) \neq 0 . \]
En conclusion :
\[ \forall x \in E, x \neq 0, \quad N'(x) = N(x) N' \left ( \dfrac{x}{N(x)} \right ) \geq b N(x) , \]
et on a exhibé $a, b$ non nuls tels que $\forall x \in E, \quad b N(x) \leq N'(x) \leq aN(x)$. 
En d'autres termes, en dimension finie, toutes les normes sont équivalentes à la norme infinie (et à fortiori entre elles).
\end{proof}

Ce théorème admet un corollaire important.

\begin{coro}
Une application linéaire d'un espace vectoriel normé $E$ de dimension finie dans un espace vectoriel normé $F$ quelconque est continue.
\end{coro}

\begin{proof}
Il suffit de munir $E$ de la norme infinie $\norminf{\cdot}$. 
En notant $(e_1, \dots, e_n)$ la base de $E$ et $a= \sum_{i=1}^n \norm{f(e_i)}$ on a :
\[ \forall x \in E, \quad  \norm{f(x)} \leq \norminf{x} \cdot a , \]
ce qui permet de déduire la continuité en $0$ de $f$ et donc sa continuité globale.
\end{proof}

Ce théorème nous fournit aussi une démonstration du lemme suivant.

\begin{lem}
Si $E$ est de dimension finie, alors toute partie fermée et bornée de $E$ est compacte.
\end{lem}

\begin{proof}
Si $E$ est de dimension finie $n$, le choix de la norme n'influe par sur la topologie induite (cf. théorème précédent), 
on choisit de munir $E$ de la norme $\norminf{\cdot}$.
On munit aussi $E$ d'une base $(e_1, \dots, e_n)$. à laquelle on associe la base $(e_1^{\ast}, \dots, e_n^{\ast})$ du dual de $E$.
Soit $(x_n)_\N$ une suite d'éléments d'une partie $F$ fermée et bornée de $E$.
Alors, pour tout $i \in \set{1, \dots, n}$, les suites $(e_i^{\ast}(x_n))_\N \subset (\K, \abs{\cdot})$ sont bornées, 
donc admettent une sous-suite convergente. 
En notant $\varphi_i$ l'extractrice correspondante à la suite $(e_i^{\ast}(x_n))_\N$, 
et en posant $\varphi = \varphi_1 \circ \dots \circ \varphi_n$ 
on obtient une extractrice pour laquelle $(x_{\varphi(n)})_\N$ est convergente (dans $F$ par fermeture). 
Donc une partie fermée et bornée de $E$ est bien compacte.
\end{proof}



\begin{thm}[de Riesz]
On suppose $E$ muni d'une norme $\norm{\cdot}$. 
Alors $E$ est de dimension finie si, et seulement si, la boule unité fermée de $E$ est compacte.
\end{thm}

\begin{proof}
Supposons que $E$ soit de dimension finie, comme $B_f(0,1)$ est fermée bornée, elle est compacte. La condition est bien nécessaire.

Réciproquement, supposons que $E$ soit de dimension infinie, 
on va montrer que $B_f(0,1)$ n'est pas précompacte et donc pas compacte. 
Par l'absurde, supposons qu'il existe $x_1, \dots, x_n$ tels que :
\[ B_f(0,1) \subset \bigcup_{i=1}^n B(x_i,1) . \]
Notons $F=\vect(x_1,\dots,x_n)$. Comme $E$ est de dimension infinie, il existe $x \in E \setminus F$.
Notons $a= d(x,x_1) < \infty$ et $\Gamma = \set{y \in F ; \norm{x-y} \leq a}$. 
L'ensemble $\Gamma$ est fermé comme image réciproque du fermé $]-\infty, a]$ par la fonction continue :
\[ \begin{array}{c c c c}
\varphi : & F & \to & \R \\
& y & \mapsto & \norm{y-x}
\end{array} \]
De plus, c'est un ensemble borné, en effet pour $y, y' \in \Gamma, \norm{y-y'} \leq 2a$.
Donc $\Gamma$ est compact, ainsi la fonction $\varphi$ atteint sa borne inférieure sur $\Gamma \subset F$ et il existe donc 
$y \in F$ tel que $\norm{y-x} = d(x,F)$. \\
Posons $x_0 = \dfrac{x-y}{\norm{x-y}}$, alors $d(x_0, F) \leq \norm{x_0-0} = 1$ (car $0 \in F$). Par ailleurs :
\[ \forall z \in F, \quad \norm{x_0-z}= \dfrac1{\norm{x-y}} \norm{x-(y+\norm{x-y}z)} \geq \dfrac{d(x,F)}{\norm{x-y}} = 1, \]
et finalement, $d(x,F)=1$.
Or $x_0 \in B_f(0,1)$, donc il existe $i$ tel que $x_0 \in B(x_i, 1)$ de sorte que $d(x_0,x_i) < 1$ ce qui est absurde par définition de
$d(x,F)=1$. Ainsi, on a démontré que le condition était suffisante.

\end{proof}

\end{document}